\answer{ex:01:1}{(a)~$8\cdot10^{10}$ \nt{GLUT}: $10$ poblaciones de la Tierra; de difusión $\approx1.9\cdot10^6$: una ciudad del tamaño de Viena. (b)~$\approx8\cdot10^{14}$ terminales, de las cuales $\approx3\cdot10^{11}$ son de difusión, fracción $\approx4\cdot10^{-4}$: $\approx16$ veces mayor que la fracción de esas neuronas ($2.2\cdot10^{-5}$).}
\answer{ex:01:2}{(a)~$\tau_c\ln(c_0/c_*)\approx4.6\ms$. (b)~La línea queda libre si $T\ge\tau_c\ln101$: hasta $\approx22$ liberaciones por segundo en lugar de $\approx220$.}
\answer{ex:01:3}{$V(s_k)=\gamma^{K-1-k}r$; $\delta_{\text{lámpara}}=\gamma^Kr=re^{-T/\tau_\gamma}\approx0.60\,r$ para cualquier $\Delta$, $\tau_\gamma\approx1.95\,\text{s}$.}
\answer{ex:01:4}{$n^*=93$, $47$, $25$, $12$: $n^*\approx5/\alpha$ para $\alpha$ pequeña.}
\answer{ex:01:5}{\nt{GLUT}: $1\to10\to10^4\to10^7$: $\approx10^7$ neuronas, $4$ eslabones, $4\ms$. \nt{DOPA}: $1\to10\to3.2\cdot10^5$: $\approx3\cdot10^5$ neuronas, $3$ eslabones, $30$ veces menos; el mismo orden que el número de neuronas \nt{DOPA}.}
\answer{ex:01:6}{$\lambda=f_EJ_E-f_IJ_I$, de donde $J_I=37.5$; razón de corrientes $f_IJ_I/f_EJ_E=0.94$; avalancha ($\lambda\ge1$) al debilitar la inhibición solo un $6.7\%$.}
\answer{ex:01:7}{$n\ge57$ presentaciones por retardo en cada estado. La misma caída la daría, por ejemplo, una imprecisión de la estimación interna del tiempo que crece con $T$.}
